\section{Introduction}\label{sec:introduction}
A shared physics-based controller must reproduce actions on bodies with
different proportions and dynamics. Changing limb lengths, collision
geometry, and mass distribution changes both the spatial realization of a
motion and the actuation needed to track it. The learning problem therefore
has two coupled axes: variation in motion and variation in body. Good
tracking on new motions for a familiar body does not establish transfer to
a new body, and a high average over bodies can conceal severe failures.

Large motion resources, shape-conditioned reference generation, and scalable
imitation provide the building blocks~\cite{amass,humanml3d,humos,phc,protomotions}.
Body-conditioned control is established~\cite{won2019}; PHC also supports
body variation, although its quantitative evaluations use the mean SMPL
body~\cite{phc}. Our question is how reliably one controller learns and
transfers across a large corpus of explicitly paired motions and bodies.
We address it through three research questions:
\begin{enumerate}
\item \textbf{RQ1: Can one shared controller learn to track diverse motions
across different bodies?} We distinguish held-out-motion performance on
selected bodies from coverage of every trained body.
\item \textbf{RQ2: How well does that capability generalise to unseen body
shapes, and where does it fail?} We evaluate inside- and outside-range bodies,
retain every shape, and examine checkpoint sensitivity.
\item \textbf{RQ3: Which design choices support that capability, and what
trade-offs do they introduce?} We compare temporal architectures and audit
reference refinement, while identifying the missing conditioning control.
\end{enumerate}

Our contribution is a paired corpus and learning workflow that preserves
clip/body identity from HUMOS generation through simulation and replay,
together with an evaluation that separates motion exposure from body
exposure. The corpus contains 20,951 clip identities and 128 body
configurations. A shared temporal slot/type-attention controller trains
directly on 18,855 clips across those bodies using streamed motion residency.
We inherit the simulator integration, PPO machinery, and motion-tracking
infrastructure from ProtoMotions; we do not claim these or morphology
conditioning itself as new. HumanML3D captions are metadata, not policy
inputs: this is reference-conditioned imitation, not text-conditioned control.

The original checkpoint achieves 99.14\% success on 1,048 independent test
clips, each paired with one selected trained body. On 233 training-split
clips across 128 unseen inside-range bodies and 128 unseen outside-range
bodies, success is 97.28\% and 97.15\%. Yet the worst bodies succeed on only
3.0\% and 13.3\% of those clips. These findings support the central claim:
\emph{one shared controller learns morphology-specific tracking with strong
average transfer to unseen shapes, but uneven robustness across individual
bodies}. The design evidence explains practical trade-offs; it does not yet
isolate the benefit of explicit morphology inputs or establish uniform
held-out-motion performance across all bodies.
